%%%%%%%%%%%%Outline%%%%%%%%%%%%%%%%%%

%%%
% message<implementation is a small kernel core, a compile-time recognizer, and operation modules>
% we implement \tool in three parts
% the kernel core is a one-time change to the verifier and JIT
% the recognizer rewrites idioms in userspace before load
% each operation ships as one or more loadable kernel modules

%%% Section 6.1 Kernel Core

%%%
% message<the kernel core is a small one-time change>
% the kernel core adds 929 lines to the Linux kernel
% the table breaks the total down by component

%%%
% message<instruction encoding: a call plus a sidecar marked via src_reg>
% each extended instruction is a pair of instructions, a call and a sidecar
% both reuse existing eBPF opcodes, a reserved src_reg value marks the pair
% the call identifies the descriptor
% the sidecar carries a packed 52-bit payload with the operands
% the low four payload bits must form a valid eBPF register number
% the recognizer re-encodes payloads whose low bits exceed the largest register
% the kernel reverses the re-encoding when reading the sidecar

%%%
% message<verifier: a lowering stage and a restore stage around the existing analysis>
% the verifier gains a lowering stage and a restore stage around its analysis
% the lowering stage walks backward and substitutes each extended instruction's proof sequence
% the stage records each original pair for restoration
% the verifier then runs unchanged on the lowered program
% the restore stage reinstates the extended instructions after success
% with proof-sequence validation and bookkeeping the verifier changes total 487 lines

%%%
% message<JIT: per-architecture dispatch, fallback resolved before the JIT>
% the x86-64 and ARM64 JITs each gain a dispatch case for extended instructions
% the JIT looks up the descriptor and decodes the sidecar payload
% the JIT then calls the descriptor's native emit for the running architecture
% with no native emit the kernel rewrites the site to its proof sequence in load-time fixups
% dispatch is 97 lines on x86-64 and 121 on ARM64

%%%
% message<module registration reuses BTF and kfunc machinery>
% a loaded module registers its descriptors for lookup by the verifier and JIT
% registration reuses BTF, the kernel's type metadata, and kfuncs, kernel functions eBPF may call
% each module exports a stub kfunc as its BTF anchor
% the module lists its descriptors in a parallel array
% the kernel caches the descriptors per program
% BTF and registration are 139 lines
% the new descriptor type and helpers add 85 lines of headers

%%% Section 6.2 Compile-Time Recognizer

%%%
% message<the recognizer is an optional pre-load rewriter with two paths>
% the recognizer rewrites idioms into extended instructions before load, through two paths
% a program built without the recognizer, or unmatched, produces the same vanilla eBPF bytecode

%%%
% message<the selector path matches idioms during LLVM instruction selection>
% the first path is a selector in the LLVM eBPF backend
% the selector matches idioms during instruction selection and emits the call and sidecar
% counted over the backend's \tool commits, the selector adds 2,126 lines, half a new selection module

%%%
% message<the bytecode-recovery path re-matches idioms on compiled bytecode>
% the second path is a bytecode-recovery pass on already-compiled bytecode
% the pass covers bytecode the selector never sees or cannot match
% the implementation is a standalone userspace tool of 6,967 lines
% the pass checks that each rewrite preserves the program's semantics
% even a wrong rewrite leaves the program safe, the verifier re-checks the lowered program
% at worst a wrong rewrite changes the program's result

%%% Section 6.3 Operation Modules

%%%
% message<a module packages an operation's two forms outside the kernel core>
% each operation is packaged as one or more kernel modules, separate from the core
% a module supplies the proof-sequence routine, the native emits, and the BTF registration

%%%
% message<modules are organized by native instruction; tree counts and scope>
% each module covers one native instruction on one architecture
% a single-architecture operation contributes modules on one side, so the counts differ
% the tree is 14 x86-64 and 11 ARM64 modules, 9,765 lines, including helpers and test-only modules

%%%%%%%%%%%%Outline%%%%%%%%%%%%%%%%%%

\section{Implementation}
\label{sec:implementation}
% 实现

We implement \tool's three parts: the one-time kernel patch to the verifier and the JIT (\S\ref{subsec:kernel}), the compile-time recognizer that rewrites idioms in userspace before load (\S\ref{subsec:llvm}), and the operation modules that ship each \lang operation as one or more loadable kernel modules (\S\ref{subsec:descriptors}).
% 我们实现 \tool 的三个部分：对验证器和 JIT 的一次性内核补丁（\S\ref{subsec:kernel}），在加载前于用户空间改写习语的编译时识别器（\S\ref{subsec:llvm}），以及把每个 \lang 操作作为一个或多个可加载内核模块发布的操作模块（\S\ref{subsec:descriptors}）。

\subsection{Kernel Patch}
\label{subsec:kernel}
% 内核补丁

The kernel patch adds 929 lines to the Linux kernel, broken down by component in Table~\ref{tab:kernel-loc}.
% 内核补丁向 Linux 内核添加了 929 行，表按组件分解了总数。
All of it is glue around the existing verifier and JIT: it leaves the verifier's analysis and the JIT's translation of ordinary instructions unchanged.
% 这些代码全部是现有验证器和 JIT 外围的胶水代码：验证器的分析和 JIT 对普通指令的翻译都保持不变。
To carry operations through this unchanged pipeline, the patch represents each \lang operation in the bytecode as a pair of instructions that reuse existing eBPF opcodes, a sidecar followed by a call.
% 为了让操作经过这条未改动的流水线，补丁在字节码中把每个 \lang 操作表示为一对重用现有 eBPF 操作码的指令：一个 sidecar，后跟一个 call。
A reserved \texttt{src\_reg} value marks the pair for \tool.
% 一个保留的 src\_reg 值把这对指令标记为 \tool 所用。
Within the pair, the call identifies the \lang operation, and the sidecar carries its operands in a packed 52-bit payload.
% 在这对指令中，call 标识 \lang 操作，sidecar 用一个打包的 52 位载荷携带操作数。

\begin{table}[t]
\centering
\small
\caption{Lines the kernel patch adds by component, measured on the \tool kernel branch over mainline Linux. The total excludes 11 lines of disassembler support and tools-side header sync.}
\label{tab:kernel-loc}
\begin{tabular}{@{}lc@{}}
\toprule
Component & Lines of code \\
\midrule
Verifier (replacement, restore, validation) & 487 \\
BTF and registration & 139 \\
JIT support (ARM64) & 121 \\
JIT support (x86-64) & 97 \\
Headers & 85 \\
\midrule
Total kernel patch & 929 \\
\bottomrule
\end{tabular}
\end{table}

% \begin{table}[t]
\centering
\small
\caption{Encoding of an operation in bytecode. A reserved \texttt{src\_reg} value marks the pair for \tool, a \texttt{BPF\_CALL} for the call and a \texttt{MOV64\_K} for the sidecar. The notation [h:l] gives a field's position within the 52-bit sidecar payload.}
\label{tab:encoding}
\begin{tabular}{@{}llcp{0.50\columnwidth}@{}}
\toprule
Insn & Field & Bits & Meaning \tabularnewline
\midrule
Call & \texttt{imm} & 32 & \raggedright BTF ID of the operation's stub kfunc \tabularnewline
 & \texttt{off} & 16 & \raggedright Index into \texttt{fd\_array} for the module's BTF file descriptor \tabularnewline
\midrule
Sidecar & \texttt{dst\_reg} & [3:0] & \raggedright Destination register or a tag that selects the operation's variant \tabularnewline
 & \texttt{off} & [19:4] & \raggedright Offset or auxiliary field per operation \tabularnewline
 & \texttt{imm} & [51:20] & \raggedright Immediate value or configuration per operation \tabularnewline
\bottomrule
\end{tabular}
\end{table}

On top of this encoding, the patch adds a step before the verifier's analysis and one after it, and leaves the analysis itself unchanged.
% 在这种编码之上，补丁在验证器分析之前和之后各加一步，分析本身保持不变。
Before the analysis, it replaces each \lang operation with the proof sequence from its module, records the original pair, and validates that the proof sequence has a single entry and a single exit, as \S\ref{subsec:safety} requires.
% 在分析之前，它用来自模块的证明序列替换每个 \lang 操作，记录原始的指令对，并按 \S\ref{subsec:safety} 的要求验证证明序列只有单一入口和单一出口。
After the analysis succeeds, it restores the pairs.
% 分析成功后，它恢复这些指令对。
In the JIT, the x86-64 and ARM64 backends each gain a case that decodes the sidecar payload and calls a function from the \lang operation's module that generates its native code for the running architecture.
% 在 JIT 中，x86-64 和 ARM64 后端各增加一个分支，解码 sidecar 载荷，并调用 \lang 操作所在模块中为当前架构生成本地代码的函数。
When a \lang operation has no native code for the architecture, the kernel rewrites the operation back to its proof sequence before the JIT runs.
% 当操作没有该架构的生成本地代码的函数时，内核在 JIT 运行前将该处重写回其证明序列。

\subsection{Compile-Time Recognizer}
\label{subsec:llvm}
% 编译时识别器

The recognizer rewrites a program's idioms into \lang operations before load, through two paths.
% 识别器在加载前通过两条路径将程序的习语重写为操作。
The first, 2{,}126 lines in the LLVM eBPF backend, matches idioms during instruction selection and emits the call and sidecar in place of the multi-instruction sequence.
% 第一条是 LLVM eBPF 后端中的选择器，它在指令选择时匹配习语，并发射调用和伴随以代替分解后的字节码；它增加了 2126 行。
The second handles programs the backend never sees or cannot match, such as those compiled by a separate toolchain, with a standalone userspace tool of 6{,}967 lines that matches the same idioms on already-compiled bytecode.
% 第二条路径处理后端看不到或无法匹配的程序，例如由其他工具链编译的程序，它是一个 6{,}967 行的独立用户空间工具，在已编译的字节码上匹配相同的习语。
Although the recognizer lies outside the TCB, it checks that each rewrite preserves the program's semantics, for example that a register the idiom overwrites is dead afterward, since a wrong rewrite would still compute a wrong result.
% 虽然识别器不在 TCB 之内，它仍会检查每次改写是否保持程序语义，例如习语覆盖的寄存器在之后不再被使用，因为错误的改写仍会算出错误的结果。
A program built without the recognizer, or one whose idioms the recognizer does not match, produces the same vanilla eBPF bytecode and loads and runs as an ordinary eBPF program, so programs can adopt \tool one at a time.
% 不经识别器构建的程序，或识别器没有匹配到习语的程序，产生相同的原生 eBPF 字节码，并作为普通 eBPF 程序加载和运行，因此程序可以逐个采用 \tool。

\subsection{Operation Modules}
\label{subsec:descriptors}
% 操作模块

Each \lang operation is packaged as one or more loadable kernel modules, separate from the kernel patch, so operations can be added without rebuilding the kernel.
% 每个 \lang 操作打包为一个或多个可加载内核模块，与内核补丁分离，因此添加操作无需重新构建内核。
A module supplies a function that produces the proof sequence and one per architecture that generates the native code, and exports a stub kfunc in BTF. The call in each pair names the operation through this kfunc.
% 模块提供一个生成证明序列的函数，以及每个架构一个生成本地代码的函数，并在 BTF 中导出一个桩 kfunc；每对指令中的 call 通过这个 kfunc 指明操作。

The module tree holds 14 modules on x86-64 and 11 on ARM64, 9{,}765 lines in all, including shared helper headers, test-only modules, and operations beyond \lang's seven (Table~\ref{tab:descriptors}).
% 模块树在 x86-64 上有 14 个模块，在 ARM64 上有 11 个，共 9765 行，包括共享辅助头文件、仅测试模块和超出 \lang 七个的操作。
Of these lines, only the functions that generate native code affect the verifier's guarantee, since the functions that produce proof sequences emit bytecode that the verifier checks.
% 这些代码中，只有生成本地代码的函数影响验证器的保证，因为生成证明序列的函数输出的字节码会被验证器检查。
Each such function has a few dozen lines and writes one or two machine instructions, and we check that it writes exactly the instructions modeled in the Lean proofs. Because the verifier never sees what these functions write, re-verifying the program cannot replace a review of this kernel code, which, like any kernel module, must be free of ordinary kernel bugs.
% 每个这样的函数只有几十行，写出一到两条机器指令，我们检查它写出的恰好是 Lean 证明中建模的指令。由于验证器看不到这些函数写出的内容，重新验证程序不能代替对这些内核代码的审查，它们与任何内核模块一样必须没有普通的内核错误。

Adding a pattern for an existing operation changes only the userspace recognizer.
% 为现有操作添加模式只修改用户空间识别器。
A new operation also needs a kernel module with parameter checks, a function that produces its proof sequence, functions that generate its native code, and registration, as well as its Lean proof (\S\ref{sec:formal-verification}). Such a module is small. The ARM64 bit-extract module, for example, has 76 lines.
% 新操作还需要一个包含参数检查、生成证明序列的函数、生成本地代码的函数和注册代码的内核模块，以及它的 Lean 证明。这样的模块很小，例如 ARM64 位提取模块只有 76 行。
Its proof joins a Lean development of 9{,}047 lines, which took about two person-days to write with help from Codex (GPT-6.1-Sol), and the artifact includes the proofs and instructions for checking them.
% 它的证明加入一个 9{,}047 行的 Lean 开发，这个开发在 Codex（GPT-6.1-Sol）的帮助下花了大约两个人日写成，制品中包含这些证明及检查它们的说明。
